% !PRINTED 541
\weberpage{541}

%<<CONT
\[
\begin{aligned}
\frac{f(\omega)}{2^{\frac{1}{2} - \frac{1}{2^{\lambda+2}}}}, & \quad m = 4^\lambda m', \quad m' \equiv 1, 2 \pmod 4, \\
% [sic] The source prints m' congruent to 4; the cube-root case on p.540 and the recurrence require 3 modulo 8.
\frac{f(\omega)}{2^{\frac{1}{2} - \frac{1}{3.\,2^{\lambda+1}}}}, & \quad m = 4^\lambda m', \quad m' \equiv 4 \pmod 8, \\
\frac{f(\omega)}{\sqrt{2}}, & \quad m = 4^\lambda m', \quad m' \equiv 7 \pmod 8.
\end{aligned}
\]

\begin{center}
\textbf{§ 143.\quad \gesperrt{Partialnormen von $f(\omega)$.}}
\end{center}

Wir machen von der Grenzformel (30) des § 141 noch eine
Anwendung auf die Bestimmung des Produktes
\weq{1}{\Pi f(\omega),}
in dem sich das Produktzeichen $\Pi$ nicht über ein vollständiges
Wurzelsystem, sondern nur über die Wurzeln $\omega$ der Formenklassen
eines Geschlechts erstreckt. Wir beschränken uns dabei aber
auf den Fall, daß
\weq{2}{\Delta = -4m}
eine \gesperrt{Stammdiskriminante} ist, daß also $m$ keinen quadratischen
Teiler habe und entweder $\equiv 1$ oder $\equiv 2 \pmod 4$ ist oder, was
dasselbe ist,
\weq{3}{\Delta \equiv -4 \ \text{oder} \ \equiv 8 \pmod{16}.}
Es sei $\delta$ ein von $\Delta$ und 1 verschiedener \gesperrt{Stammteiler} von
$\Delta$ und $\chi(\delta, k)$ der diesem Stammteiler entsprechende Charakter
der Klasse $k$.

% [sic] The second delta is printed lowercase; the complementary divisor belongs to Delta, as delta delta'=Delta immediately below confirms.
Ist $\delta'$ der zu $\delta$ komplementäre Stammteiler zu $\delta$, so ist in
diesem Falle
\[
\begin{aligned}
\delta\,\delta' &= \Delta, \\
\chi(\delta, k) &= \chi(\delta', k),
\end{aligned}
\]
und wir bekommen also alle Geschlechter, wenn wir für $\delta$ die
\gesperrt{ungeraden Stammteiler} setzen, was wir hier tun wollen. Es sei
nun wie in § 141, (29), (30):
\weq{4}{\begin{aligned}
\psi_0 &= \Bigl(2, 2, \frac{m+1}{2}\Bigr), \ \Bigl(2, 0, \frac{m}{2}\Bigr), \\
\psi &= (A, 2B, C), \\
\psi' &= \psi_0\,\psi.
\end{aligned}}
